\section{Elementary positive-map examples}

The reduction map and the automorphisms of the positive semidefinite cone are
basic structural examples of positive maps. The reduction map has a simple
trace form and an explicit Choi matrix; the automorphisms preserve the order
cone and map it onto itself.

\subsection{The reduction map}

\begin{definition}[Reduction map]
    \label{def:reduction_map}
    \lean{Matrix.reductionMap}
    \leanok
    The reduction map on $\MN{D}$, for $k\in\N$, is
    \begin{align}
        T_k(X)
         & =\tr(X)\Id-k^{-1}X.
        \label{eq:channel_reduction_map}
    \end{align}
\end{definition}

\begin{theorem}[The first reduction map is positive]
    \label{thm:reduction_map_one_positive}
    \lean{Matrix.reductionMap_one_isPositiveMap}
    \leanok
    \uses{def:positive_map, def:reduction_map}
    For $D\ge 1$, the map $T_1(X)=\tr(X)\Id-X$ on $\MN{D}$ is positive.
\end{theorem}

\begin{proof}\leanok
    If $X\ge 0$, diagonalize $X$. Its eigenvalues
    $\lambda_1,\ldots,\lambda_D$ are non-negative, and each satisfies
    $\lambda_j\le \sum_{i=1}^D\lambda_i=\tr(X)$. Therefore all eigenvalues
    of $\tr(X)\Id-X$ are non-negative.
\end{proof}

\begin{theorem}[Self-duality of the reduction map]
    \label{thm:reduction_map_self_dual}
    \lean{Matrix.traceAdjointMap_reductionMap}
    \leanok
    \uses{def:reduction_map}
    The reduction map is self-dual for the trace pairing:
    \begin{align}
        \tr(T_k(\rho)X)
         & =\tr(\rho T_k(X)).
        \label{eq:channel_reduction_self_dual}
    \end{align}
\end{theorem}

\begin{proof}\leanok
    Expanding both sides of (\ref{eq:channel_reduction_self_dual}) using
    (\ref{eq:channel_reduction_map}) and bilinearity of the trace gives
    $\tr(\rho)\tr(X)-k^{-1}\tr(\rho X)$ in each case.
\end{proof}

\begin{theorem}[Choi matrix of the reduction map]
    \label{thm:choi_matrix_reduction_map}
    \lean{ChoiJamiolkowski.choiMatrix_reductionMap}
    \leanok
    \uses{def:reduction_map}
    For $D\ge 1$, the Choi matrix of the reduction map is
    \begin{align}
        \tau(T_k)
         & =D^{-1}\Id-k^{-1}\ket{\Omega}\bra{\Omega}.
        \label{eq:channel_reduction_choi}
    \end{align}
\end{theorem}

\begin{proof}\leanok
    The trace term sends the $(a,b)$ slice of
    $\ket{\Omega}\bra{\Omega}$ to its trace times the identity, giving
    the contribution $D^{-1}\Id$. The second term is the Choi matrix of
    $k^{-1}$ times the identity map, namely
    $k^{-1}\ket{\Omega}\bra{\Omega}$.
\end{proof}

\subsection{Automorphisms of the positive semidefinite cone}

Invertible conjugations, with or without transposition, do more than preserve
positivity: they are surjective on the positive semidefinite cone.

\begin{definition}[Maps preserving the positive semidefinite cone]
    \label{def:maps_psd_cone_onto}
    \lean{MapsPSDConeOnto}
    \leanok
    \uses{def:positive_map}
    A linear map $T:\MN{D}\to\MN{D}$ \emph{maps the positive semidefinite cone
        onto itself} if $T$ is positive and every positive semidefinite matrix is
    the image under $T$ of a positive semidefinite matrix. This is condition
    (1) of \cite[Proposition~3.8]{Wolf2012Quantum}.
\end{definition}

\begin{theorem}[Conjugations preserve the positive semidefinite cone]
    \label{thm:conjugation_maps_psd_cone_onto}
    \lean{unitaryConjLM_mapsPSDConeOnto,
        unitaryConjLM_comp_transpose_mapsPSDConeOnto}
    \leanok
    \uses{def:maps_psd_cone_onto}
    Let $Y\in\MN{D}$ be invertible. Then the map $X\mapsto YXY^\dagger$ and the
    map $X\mapsto YX^TY^\dagger$ each map the positive semidefinite cone onto
    itself. This is the implication (3) $\Rightarrow$ (1) of
    \cite[Proposition~3.8]{Wolf2012Quantum}.
\end{theorem}

\begin{proof}\leanok
    Conjugation by any matrix preserves positive semidefiniteness. Moreover,
    if $W=C^\dagger C$, then
    $W^T=C^T\overline C=\overline C^\dagger\overline C\geq0$;
    hence transposition also preserves positive semidefiniteness, and both maps
    send the cone into itself. For surjectivity onto the cone, given
    a positive semidefinite $A$, set
    $W=Y^{-1}A(Y^{-1})^\dagger$, which is again positive semidefinite. Then
    $YWY^\dagger=A$ proves the conjugation case, and
    $Y(W^T)^TY^\dagger=A$ proves the transpose case, with $W$ and $W^T$
    positive semidefinite.
\end{proof}

\section{Positivity hierarchy and two-positive maps}

The following family exhibits the strict hierarchy of positivity conditions
and supplies the threshold statements used in the reduction criterion below.

\subsection{The map \texorpdfstring{$T_\eta$}{T eta} and strictness of the chain}

The cones of $k$-positive maps form a decreasing chain from the completely
positive maps, at the top amplification dimension, down to the merely positive
maps. This subsection shows that every inclusion between consecutive
dimensions, below the top one, is strict. The witnesses come from a
one-parameter family of maps whose positivity threshold is the parameter
itself, so that tuning the parameter between two integers separates the two
adjacent cones.

\begin{definition}[The family $T_\eta$]
    \label{def:t_eta}
    \lean{Matrix.tEta}
    \leanok
    For a real parameter $\eta$, the map $T_\eta$ on $\MN{D}$ is
    \begin{align}
        T_\eta(\rho) & =\tr(\rho) \Id-\eta^{-1}\rho.
        \notag
    \end{align}
    The map is defined for every real $\eta$; at $\eta=0$ it reduces to
    $\rho\mapsto\tr(\rho) \Id$, and the positivity threshold below
    assumes $\eta>0$.
\end{definition}

\begin{theorem}[Choi matrix of $T_\eta$]
    \label{thm:choi_matrix_t_eta}
    \lean{ChoiJamiolkowski.choiMatrix_tEta}
    \leanok
    \uses{def:t_eta, thm:choi_matrix_reduction_map}
    For $D\ge 1$, the Choi matrix of $T_\eta$ is
    \begin{align}
        \tau_\eta & =D^{-1}\Id-\eta^{-1}\ket{\Omega}\bra{\Omega}.
        \notag
    \end{align}
\end{theorem}

\begin{proof}
    \leanok
    The Choi matrix is linear in the map, so it splits over the two terms of
    $T_\eta$. The trace term contributes
    $\tau(\rho\mapsto\tr(\rho) \Id)=D^{-1}\Id$, since each elementary
    slice of $\ket{\Omega}\bra{\Omega}$ contributes its trace times the identity.
    The identity map contributes
    $\tau(\id)=\ket{\Omega}\bra{\Omega}$. Subtracting $\eta^{-1}$ times the
    second from the first gives
    $\tau_\eta=D^{-1}\Id-\eta^{-1}\ket{\Omega}\bra{\Omega}$.
\end{proof}

\begin{theorem}[Positivity threshold of $T_\eta$]
    \label{thm:t_eta_threshold}
    \lean{Matrix.isNPositiveMap_tEta_iff}
    \leanok
    \uses{def:t_eta, thm:choi_matrix_t_eta, thm:schmidt_rank_choi_test_iff,
        thm:maximal_overlap_schmidt_rank}
    Let $\eta>0$ and $1\le k<D$. Then $T_\eta$ is $k$-positive if and only if
    $\eta\ge k$. This is \cite[Equation~(3.11)]{Wolf2012Quantum} for $1\le k<D$;
    the top index $k=D$, where $k$-positivity is complete positivity, is
    Theorem~\ref{thm:t_eta_threshold_top}.
\end{theorem}

\begin{proof}
    \leanok
    By the Schmidt-rank Choi criterion, $k$-positivity is the nonnegativity of
    the Choi quadratic form on all vectors of Schmidt rank at most $k$. For a
    normalized such vector $\psi$, the Choi matrix
    $\tau_\eta=D^{-1}\Id-\eta^{-1}\ket{\Omega}\bra{\Omega}$ gives
    $\bra{\psi}\tau_\eta\ket{\psi}
        =D^{-1}-\eta^{-1}|\braket{\Omega}{\psi}|^2$.
    The reduced density of $\ket\Omega$ on the first factor is $\Id/D$, so by the
    maximal-overlap principle the supremum of $|\braket{\Omega}{\psi}|^2$ over
    Schmidt-rank-$k$ normalized vectors is $\|\Id/D\|_{(k)}=k/D$. The infimum of
    the quadratic form is therefore $D^{-1}-\eta^{-1}(k/D)=D^{-1}(1-k/\eta)$,
    which is non-negative precisely when $\eta\ge k$. Degree-two homogeneity
    extends the bound from normalized vectors to all vectors of Schmidt rank at
    most $k$.
\end{proof}

\begin{theorem}[Positivity threshold of $T_\eta$ at the top index]
    \label{thm:t_eta_threshold_top}
    \lean{Matrix.isNPositiveMap_tEta_card_iff}
    \leanok
    \uses{def:t_eta, thm:choi_matrix_t_eta, thm:schmidt_rank_choi_test_iff}
    Let $\eta>0$ and $D\ge 1$. Then $T_\eta$ is $D$-positive if and only if
    $\eta\ge D$.
    Together with Theorem~\ref{thm:t_eta_threshold} this is
    \cite[Equation~(3.11)]{Wolf2012Quantum} over the full range $1\le k\le D$,
    where the top index $k=D$ is complete positivity.
\end{theorem}

\begin{proof}
    \leanok
    On $\MN{D}$, $D$-positivity is complete positivity, which is
    positive semidefiniteness of the Choi operator
    $\tau_\eta=D^{-1}\Id-\eta^{-1}\ket{\Omega}\bra{\Omega}$. Its quadratic form on
    a vector $\psi$ is $D^{-1}\|\psi\|^2-\eta^{-1}|\braket{\Omega}{\psi}|^2$; since
    $\ket\Omega$ is normalized, Cauchy--Schwarz gives
    $|\braket{\Omega}{\psi}|^2\le\|\psi\|^2$ with equality at $\psi=\Omega$, so the
    form is non-negative for all $\psi$ precisely when $\eta\ge D$.
\end{proof}

\begin{corollary}[Strictness of the positivity chain]
    \label{cor:positivity_chain_strict}
    \lean{Matrix.exists_isNPositiveMap_not_succ}
    \leanok
    \uses{thm:t_eta_threshold}
    Let $1\le k$ with $k+1<D$. Then there is a map on $\MN{D}$ that is
    $k$-positive but not $(k+1)$-positive; the cone of $k$-positive maps strictly
    contains the cone of $(k+1)$-positive maps. Every inclusion between
    consecutive amplification dimensions below the top one is therefore strict.
\end{corollary}

\begin{proof}
    \leanok
    Take $T_\eta$ at $\eta=k$. By the threshold equivalence it is $k$-positive,
    since $\eta=k\ge k$, but not $(k+1)$-positive, since $\eta=k<k+1$.
\end{proof}

\begin{theorem}[Kraus-form Kadison--Schwarz as a special case]
    \lean{kadison_schwarz_from_2positive}
    \leanok
    \uses{thm:kadison_schwarz_2positive}
    Let $K=(K_i)_{i=0}^{d-1}$ be a unital Kraus family, that is,
    $\sum_i K_i K_i^\dagger=\Id$. Then for all $X\in\MN{D}$,
    $\mathcal{K}(X^\dagger X)\ge
        \mathcal{K}(X)^\dagger\mathcal{K}(X)$,
    where $\mathcal{K}(Y)=\sum_i K_i Y K_i^\dagger$.
\end{theorem}

\section{The reduction criterion}

The reduction map and the family $T_\eta$ are the same map under two
parametrizations: the natural-number parameter $k$ and the real parameter
$\eta$ agree at $\eta=k$. Identifying them lets the positivity threshold and
the Choi formula carry over to the reduction map without restating them. Wolf
attaches to this map the entanglement witness
$W_n=D^{-1}\Id-n^{-1}\ket{\Omega}\bra{\Omega}$ and the operator inequality
\cite[Equation~(3.18)]{Wolf2012Quantum} that constrains the reduced densities
of a bipartite state whose Schmidt number is at most $n$.

\begin{theorem}[The reduction map is the family $T_\eta$ at integer parameter]
    \label{thm:reduction_map_eq_t_eta}
    \lean{Matrix.reductionMap_eq_tEta}
    \leanok
    \uses{def:reduction_map, def:t_eta}
    For every $k\in\N$ the reduction map $T_k$ equals the family map $T_\eta$ at
    $\eta=k$: $T_k=T_\eta\big|_{\eta=k}$.
\end{theorem}

\begin{proof}
    \leanok
    The two maps differ only in the scalar coefficient of $X$, namely $k^{-1}$
    against $\eta^{-1}$ at $\eta=k$. The natural-number and real casts of $k$
    into the complex scalars coincide, so the coefficients agree and the maps
    are equal.
\end{proof}

\begin{theorem}[The reduction map $T_k$ is $k$-positive]
    \label{thm:reduction_map_is_n_positive}
    \lean{Matrix.reductionMap_isNPositiveMap}
    \leanok
    \uses{def:reduction_map, thm:reduction_map_eq_t_eta, thm:t_eta_threshold,
        thm:t_eta_threshold_top}
    For $1\le k\le D$ the reduction map $T_k$ on $\MN{D}$ is $k$-positive.
\end{theorem}

\begin{proof}
    \leanok
    The reduction map equals the family map $T_\eta$ at $\eta=k$, and $T_\eta$ is
    $k$-positive exactly when $\eta\ge k$; the threshold is met with equality at
    $\eta=k$. For $k<D$ this is the Schmidt-rank threshold, and the top index
    $k=D$ is the completely positive endpoint supplied by the dimension-indexed
    threshold.
\end{proof}

\begin{theorem}[The reduction map $T_k$ is not $(k+1)$-positive]
    \label{thm:reduction_map_not_n_positive_succ}
    \lean{Matrix.reductionMap_not_isNPositiveMap_succ}
    \leanok
    \uses{def:reduction_map, thm:reduction_map_eq_t_eta, thm:t_eta_threshold,
        thm:t_eta_threshold_top}
    For $1\le k<D$ the reduction map $T_k$ on $\MN{D}$ is not $(k+1)$-positive.
\end{theorem}

\begin{proof}
    \leanok
    The reduction map equals $T_\eta$ at $\eta=k$. For $k+1<D$ the threshold
    gives that $T_\eta$ is $(k+1)$-positive iff $\eta\ge k+1$, so $T_k$ would
    require $k\ge k+1$, which is false. For $k+1=D$ the dimension-indexed
    threshold gives that $T_\eta$ is $D$-positive iff $\eta\ge D$, so $T_k$ would
    require $k\ge D$, contradicting $k<D$.
\end{proof}

\begin{remark}
    Together with the $k$-positivity of $T_k$
    (\ref{thm:reduction_map_is_n_positive}), this shows that $T_k$ is
    \emph{exactly} $k$-positive for $1\le k<D$: every inclusion in Wolf's chain
    $T_{\mathrm{cp}}=T_D\subsetneq\cdots\subsetneq T_1$ is strict.
\end{remark}

\begin{definition}[Reduction witness]
    \label{def:reduction_witness}
    \lean{Matrix.reductionWitness}
    \leanok
    The reduction witness on $\MN{D}\otimes\MN{D}$, for $n\in\N$, is
    \begin{align}
        W_n & =D^{-1}\Id-n^{-1}\ket{\Omega}\bra{\Omega}.
        \notag
    \end{align}
\end{definition}

\begin{theorem}[The Choi matrix of the reduction map is the witness]
    \label{thm:choi_matrix_reduction_witness}
    \lean{ChoiJamiolkowski.choiMatrix_reductionMap_eq_reductionWitness}
    \leanok
    \uses{def:reduction_map, def:reduction_witness, thm:reduction_map_eq_t_eta,
        thm:choi_matrix_t_eta}
    For $D\ge 1$ the Choi matrix of $T_n$ is the reduction witness:
    \begin{align}
        \tau(T_n) & =W_n=D^{-1}\Id-n^{-1}\ket{\Omega}\bra{\Omega}.
        \notag
    \end{align}
\end{theorem}

\begin{proof}
    \leanok
    The reduction map is $T_\eta$ at $\eta=n$, whose Choi matrix is
    $D^{-1}\Id-n^{-1}\ket{\Omega}\bra{\Omega}$, which is exactly $W_n$.
\end{proof}

\begin{theorem}[Reduction-criterion identity]
    \label{thm:tensor_map_id_t_eta}
    \lean{Matrix.tensorMapId_tEta_eq}
    \leanok
    \uses{def:t_eta}
    Applying $T_\eta$ to the first factor of a bipartite matrix $\rho$ gives
    \begin{align}
        (T_\eta\otimes\id)(\rho) & =\Id\otimes\rho_2-\eta^{-1}\rho,
        \notag
    \end{align}
    where $\rho_2$ is the reduced density of $\rho$ on the second factor.
\end{theorem}

\begin{proof}
    \leanok
    Fix block indices $i_2,j_2$ on the second factor and write $\rho^{(i_2,j_2)}$
    for the $D\times D$ slice $\rho^{(i_2,j_2)}_{i_1 j_1}=\rho_{(i_1,i_2),(j_1,j_2)}$
    on the first factor. Its trace is the corresponding entry of $\rho_2$:
    \begin{align}
        \tr(\rho^{(i_2,j_2)})
         & =\sum_{i_1}\rho_{(i_1,i_2),(i_1,j_2)}
        =(\rho_2)_{i_2 j_2}.
        \notag
    \end{align}
    Applying $T_\eta$ entry by entry gives
    \begin{align}
        ((T_\eta\otimes\id)(\rho))_{(i_1,i_2),(j_1,j_2)}
         & =\tr(\rho^{(i_2,j_2)}) \delta_{i_1 j_1}
        -\eta^{-1}\rho_{(i_1,i_2),(j_1,j_2)}
        \notag                                     \\
         & =(\rho_2)_{i_2 j_2} \delta_{i_1 j_1}
        -\eta^{-1}\rho_{(i_1,i_2),(j_1,j_2)},
        \notag
    \end{align}
    and the first term is exactly the $(i_1,i_2),(j_1,j_2)$ entry of
    $\Id\otimes\rho_2$.
\end{proof}

\begin{theorem}[Operator-implication step of the reduction criterion]
    \label{thm:reduction_criterion}
    \lean{Matrix.reductionCriterion_left, Matrix.reductionCriterion_right}
    \leanok
    \uses{def:t_eta, thm:tensor_map_id_t_eta}
    Let $n\ge 1$ and let $\rho$ be a bipartite matrix, with $\rho^{F}$ its image
    under the factor swap $\rho^{F}_{(i_2,i_1),(j_2,j_1)}=\rho_{(i_1,i_2),(j_1,j_2)}$.
    If $(T_n\otimes\id)(\rho)\ge0$ then
    $n\,\Id\otimes\rho_2\ge\rho$, and if
    $(T_n\otimes\id)(\rho^{F})\ge0$ then
    $n\,\rho_1\otimes\Id\ge\rho$, where $\rho_1$ and $\rho_2$ are the reduced
    densities of $\rho$ on the first and second factors. Because the factor
    swap is a unitary reindexing,
    $(\id\otimes T_n)(\rho)\ge0
        \iff(T_n\otimes\id)(\rho^{F})\ge0$, so the second hypothesis is Wolf's
    symmetric condition
    $(\id\otimes T_n)(\rho)\ge0$ expressed through the swap.

    This is the second step of Wolf's two-step argument for
    \cite[Equation~(3.18)]{Wolf2012Quantum}. The source states the inequalities
    under the premise that $\rho$ has Schmidt number at most $n$, and reaches
    $(T_n\otimes\id)(\rho)\ge0$ from that premise by the $n$-positivity of
    $T_n$; that first step is taken here as the hypothesis rather than derived.
    The first step and the composed full criterion with the Schmidt-number
    premise are
    Theorem~\ref{thm:reduction_criterion_of_schmidt_number}.
\end{theorem}

\begin{proof}
    \leanok
    By the reduction-criterion identity, $(T_n\otimes\id)(\rho)$ equals
    $\Id\otimes\rho_2-n^{-1}\rho$. Scaling a positive semidefinite operator by
    $n>0$ preserves positivity, so $n\,\Id\otimes\rho_2-\rho\ge0$, which is the
    first inequality. For the second, the factor swap $F$ satisfies $F^2=\id$
    and is a unitary reindexing, so it preserves positive semidefiniteness and
    the order, and it exchanges the reduced densities and the two tensor slots:
    $(\rho^{F})_2=\rho_1$ and
    $(\Id\otimes(\rho^{F})_2-n^{-1}\rho^{F})^{F}
        =\rho_1\otimes\Id-n^{-1}\rho$.
    Applying the first form to $\rho^{F}$ gives
    $n\,\Id\otimes(\rho^{F})_2-\rho^{F}\ge0$; conjugating by $F$ turns this into
    $n\,\rho_1\otimes\Id-\rho\ge0$.
\end{proof}

\section{The Breuer--Hall map}

A second concrete positive map of \cite[Chapter~3, Example~3.1]{Wolf2012Quantum}
is built from an antisymmetric contraction $U$ on $\C^{D}$, a matrix with
$U^{\mathsf T}=-U$ and $U^{\dagger}U\le\Id$. The associated Breuer--Hall map
\cite[Equation~(3.19)]{Wolf2012Quantum} is
$T_{\mathrm{BH}}(X)=\tr(X) \Id-X-U\,X^{\mathsf T}U^{\dagger}$.
It is positive; whether it is decomposable, and the threshold beyond which it
fails to be completely positive, are not treated here.

\begin{definition}[Breuer--Hall map]
    \label{def:breuer_hall_map}
    \lean{Matrix.breuerHallMap}
    \leanok
    For a matrix $U$ on $\MN{D}$, the Breuer--Hall map on $\MN{D}$ is
    \begin{align}
        T_{\mathrm{BH}}(X) & =\tr(X) \Id-X-U\,X^{\mathsf T}U^{\dagger}.
        \notag
    \end{align}
\end{definition}

\begin{lemma}[The Breuer--Hall map on rank-one operators]
    \label{lem:breuer_hall_rankone}
    \lean{Matrix.breuerHallMap_vecMulVec_posSemidef}
    \leanok
    \uses{def:breuer_hall_map}
    If $U$ is antisymmetric, $U^{\mathsf T}=-U$, and a contraction,
    $U^{\dagger}U\le\Id$, then for every vector $v$ the image
    $T_{\mathrm{BH}}(\ket{v}\bra{v})$ of the rank-one operator $\ket{v}\bra{v}$ is
    positive semidefinite.
\end{lemma}

\begin{proof}
    \leanok
    \uses{def:breuer_hall_map}
    On $\ket{v}\bra{v}$ the image is
    \begin{align}
        T_{\mathrm{BH}}(\ket{v}\bra{v})
         & =\lVert v\rVert^{2} \Id-\ket{v}\bra{v}
        -\ket{U\overline{v}}\bra{U\overline{v}},
        \notag
    \end{align}
    with $\overline{v}$ the entrywise conjugate of $v$. The two vectors $v$ and
    $U\overline{v}$ are orthogonal: antisymmetry of $U$ gives
    \begin{align}
        \braket{v}{U\overline{v}}
         & =\overline{v}^{\,\mathsf T}U\,\overline{v}
        =\overline{v}^{\,\mathsf T}(-U^{\mathsf T})\overline{v}
        =-\braket{v}{U\overline{v}},
        \notag
    \end{align}
    hence $\braket{v}{U\overline{v}}=0$. The contraction bound gives
    $\lVert U\overline{v}\rVert\le\lVert v\rVert$. For any
    $w$ the quadratic form is
    $\lVert v\rVert^{2}\lVert w\rVert^{2}
        -\lvert\braket{v}{w}\rvert^{2}
        -\lvert\braket{U\overline{v}}{w}\rvert^{2}$,
    which is non-negative because the two-vector Bessel inequality for the
    orthogonal pair $v$, $U\overline{v}$ with
    $\lVert U\overline{v}\rVert\le\lVert v\rVert$ gives
    $\lvert\braket{v}{w}\rvert^{2} +\lvert\braket{U\overline{v}}{w}\rvert^{2}
        \le\lVert v\rVert^{2}\lVert w\rVert^{2}$.
\end{proof}

\begin{theorem}[The Breuer--Hall map is positive]
    \label{thm:breuer_hall_positive}
    \lean{Matrix.breuerHallMap_isPositiveMap}
    \leanok
    \uses{def:breuer_hall_map}
    If $U$ is antisymmetric, $U^{\mathsf T}=-U$, and a contraction,
    $U^{\dagger}U\le\Id$, then the Breuer--Hall map $T_{\mathrm{BH}}$ sends every
    positive semidefinite matrix to a positive semidefinite matrix.
\end{theorem}

\begin{proof}
    \leanok
    \uses{lem:breuer_hall_rankone}
    A positive semidefinite matrix is a sum of rank-one operators
    $\ket{v}\bra{v}$. By Lemma~\ref{lem:breuer_hall_rankone} the Breuer--Hall map
    sends each such operator to a positive semidefinite matrix, and a sum of
    positive semidefinite matrices is positive semidefinite, so $T_{\mathrm{BH}}$
    is positive.
\end{proof}
